\addtocontents{toc}{\protect}
\section{\texorpdfstring{\MakeUppercase{future work and perspectives}}{future work and perspectives}}
\label{sec:chapter-future}

\counterwithin{figure}{section}
\counterwithin{table}{section}


\subsection{Updates to McFACTS}
\texttt{McFACTS} is designed to be public, modular and adaptable.
The content presented in this dissertation was performed using \code{v0.3.0}, but we have already implemented stars into the disk (Nathaniel et al. 2026 in prep.) in \code{v0.4.0}.
Other future near-term goals include:
(i) add objects other than BH and main-sequence stars to the disk (neutron stars, white dwarfs)
(ii) consider the EM significance of encounters between objects (including TDEs, partial TDEs, GRBs as well as disk-crossing events),
(iii) treat the final time step of mergers with increasing sophistication, including calculating details of eccentricity evolution, recoil kicks from numerical relativity calculations (Ray et al. 2026 in prep.) and disk re-capture correctly,
(iv) allow multiple AGN episodes to occur in a given nucleus, with appropriate dynamical relaxation, (v) predict the rate and magnitude of optical/ultraviolet/X-ray flaring due to explosive events \citep[mergers, disk-crossings, SNe, TDEs,][]{McPike+26-mcfactsIV},
(vi) develop and add new models of AGN disks to reflect recent simulations \citep{Hopkins+24}, and
(vii) add more sophisticated treatments for binary capture \citep{Whitehead+25b} and stalling of gas hardening (Postiglione et al. 2026 in prep.).

\subsection{Updates to Ocotillo and New Models}


The embedded SNe light curve we produced was marginally visible against a bright AGN disk and the color evolution would likely be lost in the noise of typical AGN variability. 
However, we did not capture the peak of the light curve due to domain size limitations with our uniform grid.
Future studies should move to a code that supports adaptive mesh refinement or static mesh refinement at a minimum since we are approaching the limits of resolution for three-dimensional uniform grids in terms of computational expense and data storage requirements.
Additionally, we will likely need to move to a global disk model that will enable us to model the full disk emission.
Codes like \citep[\code{Athena++}][]{Stone+20} may be appropriate for this purpose.

Ocotillo will need an upgrade to calculate in three dimensions for a global disk in spherical coordinates that spans the full meridional domain rather than one with boundaries just above and below the disk.
We will also expand its treatment for the gas composition to achieve ionization states for a broader range of atomic species that more accurately reflect the solar and super-solar metallicities found in AGN \citep{HuangLinShields23}, requiring additional opacities starting with helium.


We also plan to model the hydrodynamics and radiative transfer of a recoiling BBH merger remnant.
In a black hole merger, some mass is lost and converted into gravitational waves, dissipating energy and angular momentum; a loss that is accompanied by a decrease of the Hill radius of the remnant in comparison with the combined Hill radius of the progenitors.
The dynamical consequences are manifold.
First, gas within the Hill sphere that is orbiting too fast for the new lower mass moves outward, self-colliding \citep{McKernan+19}.
Second, gas no longer bound by the new, smaller Hill sphere shocks with AGN gas in orbit around the central mass.
Third, a gravitational wave kick at merger causes ram-pressure stripping of the original Hill sphere gas as it collides with a comparable mass of disk gas.
Flares from disk-remnant interactions have been studied analytically \citep{McKernan+19}, showing that whether the flare is prompt or delayed depends on whether the Hill sphere radius is larger than (prompt ultraviolet and optical flare) or smaller than (delayed and diminished with drawn out optical flare) the disk scale height.
Our proposed work will test these analytical results and produce theoretical observables for comparison with follow-up campaigns to future gravitational wave detections, shedding light onto a problem that has received little numerical attention \citep{Corrales+10}.
To model these disk-remnant interactions, we will also need to move to a code with static or adaptive mesh refinement for a global model.
We will represent the recoiling black hole by giving a velocity kick to the portion of
the disk gas within its Hill sphere directed along a chosen trajectory while varying the direction and kick velocity.
This construction eliminates the need for a domain large enough to contain the remnant’s
entire orbital path, and focuses computational resources on the disk-remnant interaction that occurs
over a short timescale in a localized region of the disk.

